%%%%%%%%%%%%%%%%%%%%%%%%%%%%%%%%%%%%%%%%%%%%%%%%%%%%%%%%%%%%%%%%%%%%%%
%%                     Omitted Process
%%%%%%%%%%%%%%%%%%%%%%%%%%%%%%%%%%%%%%%%%%%%%%%%%%%%%%%%%%%%%%%%%%%%%%
%\color{blue}
\subsection{Glyph: \glyph{Omitted process}}\label{sec:omitted}

Omitted processes are processes that are known to exist, but are omitted from the diagram for the sake of clarity or parsimony. A single \glyph{omitted process} can represent any number of actual processes. The omitted process is different from a module.

\begin{glyphDescription}
 \item[SBO]\mbox{}\\ To be determined
 \item[origin]\mbox{}\\ One or several \glyph{consumption} arcs (section \ref{sec:consumption}) or one or several \glyph{production} arcs (section~\ref{sec:production}).
 \item[target]\mbox{}\\ One or several \glyph{production} arcs(section \ref{sec:production}).
 \item[node]\mbox{}\\ Omitted processes are represented as a transition containing a double slanted NW-SE line separated by an empty space.
 \end{glyphDescription}

\begin{figure}[H]
  \centering
  \includegraphics[scale = 0.5]{images/omitted}
  \caption{The \PD glyph for \glyph{omitted}.}
  \label{fig:omitted}
\end{figure}



